% Quiz 3, Part C: Stochastic Processes. Included by quiz03.tex (solutions
% hidden) and quiz03_solutions.tex (solutions shown). Problems C1--C22, grouped
% by the battery of the Stochastic Processes lecture they practise; within a
% group they get harder. Every number was computed in python3
% (scratchpad/quizC.py, sympy/fractions/numpy); values are quoted in comments.
\providecommand{\Var}{\operatorname{Var}}
\providecommand{\Cov}{\operatorname{Cov}}
\providecommand{\PP}{\mathbf{P}}
\providecommand{\st}[1]{\mathsf{#1}}
\providecommand{\Z}{\mathbb{Z}}

\noindent Notation as in the Stochastic Processes lecture: $\xi_1,\xi_2,\dots$ are the $\pm1$ steps, $S_n=\xi_1+\dots+\xi_n$, and $\Prob_i,\E_i$ refer to a start at $i$. $\PP=(p(i,j))$ is the transition matrix, row $i$ the conditional pmf of the next state given the current state $i$; $p^{(n)}(i,j)$ are the entries of $\PP^n$; $\pi$ is a stationary distribution; $T_A$ is the hitting time of $A$, $h^A_i=\Prob_i(T_A<\infty)$, $k^A_i=\E_i[T_A]$; $\mathbf N=(\mathbf I-\mathbf Q)^{-1}$ is the fundamental matrix. $\Phi$ is the $N(0,1)$ cdf; values used below:
\begin{center}\small
\begin{tabular}{@{}l*{9}{c}@{}}
\toprule
$x$ & $0.100$ & $0.500$ & $0.577$ & $0.674$ & $0.707$ & $0.905$ & $1.000$ & $1.155$ & $1.225$ \\
$\Phi(x)$ & $0.540$ & $0.691$ & $0.718$ & $0.750$ & $0.760$ & $0.817$ & $0.841$ & $0.876$ & $0.890$ \\
\bottomrule
\end{tabular}
\qquad $\Phi(1.414)=0.921$.
\end{center}

\begin{enumerate}[label=C\arabic*.,leftmargin=2.6em]
\subsection*{Processes and the random walk \hfill\normalfont\small\color{soft}Lecture batteries 1--2}
\item Three fair coin tosses, $\xi_t=\pm1$, $S_t=\xi_1+\dots+\xi_t$: the eight-row table of the lecture.
(a) Write the joint pmf of $(S_1,S_3)$ as a table and check that $S_1$ and $S_3$ are not independent.
(b) Read $\E[S_3\mid S_1=s]$ and $\E[S_3\mid S_2=s]$ off the rows, for every $s$, and explain the answers by independent increments.
(c) Compute $\E[S_3^2\mid S_2=s]$ for $s=-2,0,2$ and write $\E[S_3^2\mid S_2]$ as a random variable.
\begin{solution}
(a) Each row has probability $\tfrac18$; $(S_1,S_3)$ takes $(1,3)$ on $\st{HHH}$, $(1,1)$ on $\st{HHT},\st{HTH}$, $(1,-1)$ on $\st{HTT}$, $(-1,1)$ on $\st{THH}$, $(-1,-1)$ on $\st{THT},\st{TTH}$, $(-1,-3)$ on $\st{TTT}$:
\begin{center}
\begin{tabular}{@{}c|cccc@{}}
$p_{S_1,S_3}$ & $S_3=-3$ & $S_3=-1$ & $S_3=1$ & $S_3=3$ \\ \hline
$S_1=1$ & $0$ & $\tfrac18$ & $\tfrac28$ & $\tfrac18$ \\
$S_1=-1$ & $\tfrac18$ & $\tfrac28$ & $\tfrac18$ & $0$
\end{tabular}
\end{center}
So $\Prob(S_1=1,S_3=1)=\tfrac28\ne\Prob(S_1=1)\Prob(S_3=1)=\tfrac12\cdot\tfrac38$.
(b) Rows with $S_1=1$: $S_3=3,1,1,-1$, average $1$; with $S_1=-1$: average $-1$. Rows with $S_2=2$: $S_3=3,1$, average $2$; $S_2=0$: $S_3=1,-1,1,-1$, average $0$; $S_2=-2$: average $-2$. In each case $\E[S_3\mid S_m=s]=s$, because $S_3=S_m+(\xi_{m+1}+\dots+\xi_3)$ and the increment is independent of $S_m$ with mean $0$.
(c) $S_3^2$ on the rows with $S_2=2$: $9,1$, average $5$; $S_2=0$: $1,1,1,1$, average $1$; $S_2=-2$: $5$. So $\E[S_3^2\mid S_2=s]=s^2+1$ and $\E[S_3^2\mid S_2]=S_2^2+1$: the square of the known position plus the variance of one step.
% python3: joint counts (1,3):1,(1,1):2,(1,-1):1,(-1,1):1,(-1,-1):2,(-1,-3):1; E[S3^2|S2=2]=5, =0:1
\end{solution}

\item A walk with $p=0.55$ (a slightly favourable game), $n=100$ steps.
(a) Compute $\E[S_{100}]$, $\Var(S_{100})$ and the standard deviation.
(b) Compute $\Prob(S_{100}=10)$ and $\Prob(S_{100}=11)$.
(c) Approximate $\Prob(S_{100}>0)$ by the normal distribution (use $\{S_{100}\ge2\}$ and the midpoint $1$ as the cut).
(d) Compute $\Cov(S_{30},S_{100})$ and the correlation of $S_{30}$ and $S_{100}$.
(e) Fair walk: compute $\Prob(S_8=0)$, $\Prob(S_8=4)$ and $\Prob(S_9=0)$.
\begin{solution}
(a) $\E[S_{100}]=100(0.55-0.45)=10$; $\Var(S_{100})=4\cdot100\cdot0.55\cdot0.45=99$; standard deviation $9.95$.
(b) $S_{100}=10$ iff $H_{100}=55$: $\binom{100}{55}0.55^{55}0.45^{45}=0.080$. $S_{100}=11$ is impossible: $S_{100}$ is even.
(c) $\Prob(S_{100}\ge2)\approx\Prob\big(N(10,99)>1\big)=1-\Phi\big(\tfrac{1-10}{9.95}\big)=\Phi(0.905)=0.817$; the exact binomial sum is $0.817$.
(d) $S_{100}=S_{30}+(S_{100}-S_{30})$ with the increment independent of $S_{30}$, so $\Cov(S_{30},S_{100})=\Var(S_{30})=4\cdot30\cdot0.55\cdot0.45=29.7$; correlation $29.7/\sqrt{29.7\cdot99}=\sqrt{30/100}=0.548$.
(e) $\Prob(S_8=0)=\binom84/256=\tfrac{70}{256}=\tfrac{35}{128}$; $\Prob(S_8=4)=\binom86/256=\tfrac{28}{256}=\tfrac7{64}$; $\Prob(S_9=0)=0$ by parity.
% python3: C(8,4)/256=35/128, C(8,6)/256=7/64
% python3: P(S100=10)=0.07999; exact P(S100>0)=0.81727, normal with cut at 1: 0.81714; Cov 29.7, corr 0.5477
\end{solution}

\item Fair walk, $M_8=\max_{k\le8}S_k$.
(a) Compute $\Prob(M_8\ge2)$ by the reflection principle.
(b) Compute $\Prob(M_8=0)$.
(c) Tabulate $\Prob(M_8\ge m)$ for $m=1,\dots,8$ and deduce the pmf of $M_8$.
(d) Compute $\E[M_8]=\sum_{m\ge1}\Prob(M_8\ge m)$. What would the head-on computation require?
\begin{solution}
(a) $\Prob(M_8\ge2)=2\Prob(S_8>2)+\Prob(S_8=2)=2\Prob(S_8\ge4)+\Prob(S_8=2)=2\cdot\frac{28+8+1}{256}+\frac{56}{256}=\frac{130}{256}=\frac{65}{128}$.
(b) $\Prob(M_8=0)=\Prob(S_8=0)+\Prob(S_8=1)=\frac{70}{256}+0=\frac{35}{128}$: over a quarter of all paths never rise above their start.
(c) $\Prob(M_8\ge m)$ for $m=1,\dots,8$: $\frac{93}{128},\frac{65}{128},\frac{37}{128},\frac{23}{128},\frac{9}{128},\frac{5}{128},\frac{1}{128},\frac{1}{256}$. Differences give the pmf of $M_8$ on $0,1,\dots,8$: $\frac{70,\,56,\,56,\,28,\,28,\,8,\,8,\,1,\,1}{256}$, which is $\Prob(S_8=m)+\Prob(S_8=m+1)$ as in the lecture's corollary.
(d) $\E[M_8]=\frac{186+130+74+46+18+10+2+1}{256}=\frac{467}{256}=1.824$. Head-on: track the running maximum along each of the $256$ paths.
% python3: brute force over 256 paths agrees: pmf (70,56,56,28,28,8,8,1,1)/256, E = 467/256
\end{solution}

\item A fair walk on $\{0,1,\dots,a\}$ with a reflecting wall: from $0$ the next position is $1$ with probability one; from $0<i<a$ it moves to $i\pm1$ with probability $\tfrac12$ each; $a$ is absorbing. Let $m_i=\E_i[\text{time to reach }a]$.
(a) Write the first-step equations for $m_0,\dots,m_a$.
(b) Solve them for $a=3$.
(c) Show that $m_i=a^2-i^2$ for every $a$, and explain the agreement with the lecture's exit time from $(-a,a)$.
\begin{solution}
(a) $m_a=0$; $m_0=1+m_1$; $m_i=1+\tfrac12m_{i-1}+\tfrac12m_{i+1}$ for $0<i<a$.
(b) $a=3$: $m_3=0$, $m_2=1+\tfrac12m_1$, $m_1=1+\tfrac12m_0+\tfrac12m_2$, $m_0=1+m_1$. Substituting, $m_1=1+\tfrac12(1+m_1)+\tfrac12(1+\tfrac12m_1)$, so $\tfrac14m_1=2$: $m_1=8$, $m_0=9$, $m_2=5$.
(c) $m_i=a^2-i^2$: at $i=a$ it is $0$; at $0<i<a$, $1+\tfrac12(a^2-(i-1)^2)+\tfrac12(a^2-(i+1)^2)=a^2-i^2$; at $0$, $1+(a^2-1)=a^2$. Uniqueness as in the lecture. If $X_n$ is the unrestricted fair walk started at $i$, then $|X_n|$ is exactly this reflected walk, and $|X_n|$ reaches $a$ when $X_n$ leaves $(-a,a)$: the same random time, so the same expectation $a^2-i^2$.
% python3 sympy: a=3 -> (9,8,5,0); a=5 -> (25,24,21,16,9,0)
\end{solution}

\subsection*{Brownian motion \hfill\normalfont\small\color{soft}Lecture battery 3}
\item Let $B$ be a standard Brownian motion.
(a) Give the distribution of $B_3-B_1$ and compute $\Prob(B_3-B_1>2)$.
(b) Compute $\Prob(B_4>1\mid B_1=2)$ and $\Prob(B_4>1\mid B_1=-1)$ by splitting into increments.
(c) Write $\E[B_4\mid B_1]$ and $\E[B_4^2\mid B_1]$ as random variables.
(d) Now condition the other way. Explain by symmetry between the three increments $B_1$, $B_2-B_1$, $B_3-B_2$ why $\E[B_1\mid B_3]=B_3/3$. Then, by completing the square in the joint pdf of $(B_1,B_3)$, show that given $B_3=x$, $B_1\sim N(\tfrac x3,\tfrac23)$, and compute $\Prob(B_1>0\mid B_3=3)$.
\begin{solution}
(a) $B_3-B_1\sim N(0,2)$; $\Prob(B_3-B_1>2)=1-\Phi(2/\sqrt2)=1-\Phi(1.414)=0.079$.
(b) $B_4=B_1+(B_4-B_1)$ with $B_4-B_1\sim N(0,3)$ independent of $B_1$. Given $B_1=2$: $B_4\sim N(2,3)$, $\Prob(B_4>1)=\Prob(N(0,1)>-1/\sqrt3)=\Phi(0.577)=0.718$. Given $B_1=-1$: $B_4\sim N(-1,3)$, $\Prob=1-\Phi(2/\sqrt3)=1-\Phi(1.155)=0.124$.
(c) $\E[B_4\mid B_1]=B_1+\E[B_4-B_1]=B_1$; $\E[B_4^2\mid B_1]=B_1^2+2B_1\E[B_4-B_1]+\E[(B_4-B_1)^2]=B_1^2+3$.
(d) The three increments are i.i.d.\ $N(0,1)$ and $B_3$ is their sum, so given $B_3$ each has the same conditional mean and the three means add to $B_3$: $\E[B_1\mid B_3]=B_3/3$ (as $\E[S_5\mid S_8]=\tfrac58S_8$ in the Conditional Expectation lecture, battery 4, exercise 4). Joint pdf: $f(x,y)=\frac1{2\pi\sqrt2}e^{-x^2/2}e^{-(y-x)^2/4}$. As a function of $x$ the exponent is $-\tfrac{x^2}2-\tfrac{(y-x)^2}4=-\tfrac34\big(x-\tfrac y3\big)^2+\text{const}(y)$, so $f_{B_1\mid B_3}(x\mid y)\propto e^{-(x-y/3)^2/(2\cdot2/3)}$: $N(\tfrac y3,\tfrac23)$. Then $\Prob(B_1>0\mid B_3=3)=\Prob\big(N(1,\tfrac23)>0\big)=\Phi(1/\sqrt{2/3})=\Phi(1.225)=0.890$.
% python3: 1-Phi(sqrt2)=0.0786; Phi(1/sqrt3)=0.7181; 1-Phi(2/sqrt3)=0.1241; sympy conditional pdf integrates to 0.88966
\end{solution}

\item Use $\Cov(B_s,B_t)=\E[B_sB_t]=\min(s,t)$ throughout.
(a) Compute $\Var(B_1+B_2+B_3)$.
(b) Compute $\Var(B_2-2B_1)$ and $\Cov(B_2-2B_1,B_1)$. For which constant $c$ is $B_2-cB_1$ uncorrelated with $B_1$? Which property of Brownian motion says that for this $c$ it is in fact independent of $B_1$?
(c) Show $\operatorname{Corr}(B_s,B_t)=\sqrt{s/t}$ for $s<t$, and evaluate it for $(1,4)$ and $(1,100)$.
(d) Let $X_t=B_{t+1}-B_t$ for $t\ge0$. Show that $X_t\sim N(0,1)$ for every $t$ and that $\Cov(X_s,X_t)=\max(0,\,1-|t-s|)$. In what sense is $X$ stationary although $B$ is not?
\begin{solution}
(a) $\sum_{s,t\in\{1,2,3\}}\min(s,t)=(1+2+3)+2(1+1+2)=14$.
(b) $\Var(B_2-2B_1)=\Var(B_2)-4\Cov(B_2,B_1)+4\Var(B_1)=2-4+4=2$; $\Cov(B_2-2B_1,B_1)=1-2=-1\ne0$, so not independent. $\Cov(B_2-cB_1,B_1)=1-c=0$ iff $c=1$: the only such combination is the increment $B_2-B_1$, and property (ii) (independent increments) says that this increment is independent of $B_1=B_1-B_0$, not merely uncorrelated.
(c) $\frac{\min(s,t)}{\sqrt{st}}=\frac{s}{\sqrt{st}}=\sqrt{s/t}$; $\sqrt{1/4}=0.5$ and $\sqrt{1/100}=0.1$.
(d) $X_t$ is an increment over an interval of length $1$, so $N(0,1)$. For $s\le t$: $\Cov(X_s,X_t)=\min(s+1,t+1)-\min(s+1,t)-\min(s,t+1)+\min(s,t)=(s+1)-\min(s+1,t)-s+s=s+1-\min(s+1,t)$, which is $1-(t-s)$ if $t\le s+1$ and $0$ if $t>s+1$. Every $X_t$ has the same distribution and the covariance depends only on $t-s$: the process of unit increments looks the same at every time, while $\Var(B_t)=t$ grows.
% python3: Var=14; Var(B2-2B1)=2; Cov example s=0,t=0.5: 0.5
\end{solution}

\item The reflection principle for Brownian motion: $\Prob(\max_{s\le t}B_s\ge a)=2\Prob(B_t\ge a)$. Let $T_1$ be the first time $B$ reaches $1$.
(a) Compute $\Prob(\max_{s\le4}B_s\ge2)$.
(b) Write $\Prob(T_1\le t)$ in terms of $\Phi$ and evaluate it at $t=1,4,100$.
(c) Find the median of $T_1$ (the $t$ with $\Prob(T_1\le t)=\tfrac12$).
(d) By scaling, what is the median of $T_a$?
(e) The lecture showed $\E[T_1]=\infty$. Reconcile this with a median of about $2$.
\begin{solution}
(a) $2\Prob(B_4\ge2)=2\Prob(N(0,1)\ge1)=2(1-\Phi(1))=0.317$.
(b) $\{T_1\le t\}=\{\max_{s\le t}B_s\ge1\}$, so $\Prob(T_1\le t)=2\big(1-\Phi(1/\sqrt t)\big)$: $0.317$ at $t=1$; $2(1-\Phi(0.5))=0.617$ at $t=4$; $2(1-\Phi(0.1))=0.920$ at $t=100$.
(c) $2(1-\Phi(1/\sqrt t))=\tfrac12$ iff $\Phi(1/\sqrt t)=\tfrac34$ iff $1/\sqrt t=0.674$: median $t=1/0.674^2=2.20$.
(d) $\tfrac1a B_{a^2t}$ is a Brownian motion, and it reaches $1$ exactly when $B$ reaches $a$, at time $a^2$ times smaller; so $T_a$ has the distribution of $a^2T_1$ and median $2.20\,a^2$.
(e) Half the paths reach $1$ within $2.2$ time units, but the tail is heavy: $\Prob(T_1>t)=2\Phi(1/\sqrt t)-1\approx\sqrt{2/(\pi t)}$, and $\int^\infty t^{-1/2}\dd t=\infty$. A typical wait is short; the rare long waits are long enough to make the mean infinite.
% python3: 2(1-Phi(1))=0.3173; t=4: 0.6171; t=100: 0.9203; median 2.198
\end{solution}

\subsection*{Risk of ruin \hfill\normalfont\small\color{soft}Lecture batteries 4--5}
\item A gambler with $\pounds20$ plays $\pounds1$ rounds until she has $\pounds50$ or nothing.
(a) Fair game: compute the probability of reaching $\pounds50$ and the expected number of rounds.
(b) $p=0.48$: the same, to three decimals.
(c) She switches to $\pounds10$ rounds. Redo (a) and (b) (a walk from $2$ to $5$); verify the unfair probability by solving the four first-step equations directly rather than by the formula.
(d) State in one sentence what the stake changes and what it does not.
\begin{solution}
(a) $h_{20}=\tfrac{20}{50}=0.4$; $D_{20}=20\cdot30=600$ rounds.
(b) $r=\tfrac{0.52}{0.48}=\tfrac{13}{12}$; $h_{20}=\frac{r^{20}-1}{r^{50}-1}=0.074$; $D_{20}=\frac{20}{0.04}-\frac{50}{0.04}\cdot\frac{1-r^{20}}{1-r^{50}}=500-1250\cdot0.0737=408$ rounds.
(c) Fair: $h_2=\tfrac25=0.4$, $D_2=2\cdot3=6$ rounds. Unfair: $h_0=0$, $h_5=1$, $h_i=0.48h_{i+1}+0.52h_{i-1}$ for $i=1,\dots,4$; eliminating, $h_2=\tfrac{43200}{122461}=0.353$ (the formula gives $\frac{r^2-1}{r^5-1}$, the same), and $D_2=5.90$ rounds.
(d) In the fair game the stake changes only the duration; in the unfair game bigger stakes mean fewer rounds for the edge to compound, and the chance of reaching the target rises from $0.074$ to $0.353$.
% python3 fractions: (20,50,.5): .4, 600; (20,50,.48): .07367, 407.9; (2,5,.5): .4, 6; (2,5,.48): 43200/122461=.35277, 5.904
\end{solution}

\item Red at roulette: European wheels have one zero ($p=\tfrac{18}{37}$), American wheels two ($p=\tfrac{18}{38}$). A player starts with $\pounds100$ and bets $\pounds1$ a spin.
(a) Compute the probability of reaching $\pounds200$ on each wheel, and the ratio.
(b) Compute the expected number of spins on each wheel.
(c) European wheel, modest target $\pounds110$: compute $h_{100}$ and $\E[T]$, and check the identity $\text{mean gain}=Nh_i-i=(p-q)\,\E_i[T]$ numerically.
(d) With no target, ruin is certain on both wheels. Compute the expected number of spins until ruin on each.
\begin{solution}
(a) European: $r=\tfrac{19}{18}$, $h_{100}=\frac{r^{100}-1}{r^{200}-1}=0.00447$. American: $r=\tfrac{10}9$, $h_{100}=0.0000266$. A second zero divides the chance by about $170$.
(b) $D_{100}=\frac{100}{q-p}-\frac{200}{q-p}\cdot\frac{1-r^{100}}{1-r^{200}}$: European $q-p=\tfrac1{37}$, $D=3667$ spins; American $q-p=\tfrac2{38}$, $D=1900$ spins.
(c) $h_{100}=\frac{r^{100}-1}{r^{110}-1}=0.581$, $\E[T]=1334$ spins. Mean gain $110\cdot0.581-100=-36.1$; and $(p-q)\E[T]=-\tfrac1{37}\cdot1334=-36.1$.
(d) $\frac{i}{q-p}$: European $100\cdot37=3700$ spins; American $100\cdot19=1900$ spins.
% python3 fractions: E: h=0.004466, D=3666.9; A: h=2.656e-5, D=1899.9; E 100->110: h=0.58126, D=1334.3, gain -36.1; 3700, 1900
\end{solution}

\item Bold play. $p=0.4$, fortune $\pounds3$, target $\pounds10$.
(a) Timid play ($\pounds1$ per round): compute $h_3$ and $\E_3[T]$.
(b) Bold play: each round stake $\min(i,10-i)$, so the fortune moves from $3$ to $6$ or $0$, from $6$ to $10$ or $2$, from $2$ to $4$ or $0$, from $4$ to $8$ or $0$, from $8$ to $10$ or $6$. Draw the chain on $\{0,2,3,4,6,8,10\}$ and write the hitting system for $h_i=\Prob_i(\text{reach }10)$.
(c) Solve it and compare $h_3$ with (a).
(d) Compute the expected number of bets under bold play from $\pounds3$.
(e) Repeat (c) for the fair game $p=\tfrac12$. Why is the answer the same as for timid play there?
\begin{solution}
(a) $r=1.5$: $h_3=\frac{1.5^3-1}{1.5^{10}-1}=0.042$; $\E_3[T]=\frac{3}{0.2}-\frac{10}{0.2}\cdot\frac{1-1.5^3}{1-1.5^{10}}=12.9$ rounds.
(b) $h_0=0$, $h_{10}=1$, $h_3=ph_6$, $h_6=p+qh_2$, $h_2=ph_4$, $h_4=ph_8$, $h_8=p+qh_6$.
(c) From $h_8=p+qh_6$ and $h_6=p+qh_2=p+qp^2h_8$: $h_8=p+qp+q^2p^2h_8$, so $h_8=\frac{p(1+q)}{1-q^2p^2}=\frac{0.64}{0.9424}=\frac{400}{589}=0.679$; then $h_6=\frac{274}{589}=0.465$, $h_4=\frac{160}{589}=0.272$, $h_2=\frac{64}{589}=0.109$, and $h_3=ph_6=\frac{548}{2945}=0.186$: four and a half times the timid $0.042$.
(d) $k_0=k_{10}=0$, $k_3=1+pk_6$, $k_6=1+qk_2$, $k_2=1+pk_4$, $k_4=1+pk_8$, $k_8=1+qk_6$: $k_3=\frac{1073}{589}=1.82$ bets, against $12.9$ rounds of timid play.
(e) With $p=\tfrac12$ the system gives $h_3=\tfrac3{10}$, the timid value. A fair game preserves expected fortune whatever the stakes, so $\E[X_T]=3=10h_3$ for every strategy that stops at $0$ or $10$.
% python3 sympy: bold h = (64,548/5,160,274,400)/589 for i=2,3,4,6,8 -> h3=548/2945=0.1861; k3=1073/589=1.822; timid 0.04191, 12.90; fair bold h3=3/10
\end{solution}

\subsection*{Markov chains, paths and $n$-step transitions \hfill\normalfont\small\color{soft}Lecture batteries 6--7}
\item A taxi works three zones $\st A,\st B,\st C$. The zone where a trip ends, hence where the next trip starts, depends only on where the current trip started:
\[
  \PP=\begin{pmatrix}0.5&0.3&0.2\\0.2&0.6&0.2\\0.3&0.3&0.4\end{pmatrix}\qquad(\text{rows and columns in the order }\st A,\st B,\st C).
\]
$X_n$ is the starting zone of trip $n$.
(a) Check that $\PP$ is stochastic and say in words what row $\st B$ means.
(b) Compute $\Prob(X_1=\st C,X_2=\st B\mid X_0=\st A)$.
(c) Compute row $\st A$ of $\PP^2$ by hand and read off $\Prob(X_2=\st C\mid X_0=\st A)$. Which three two-trip sequences does that entry add up?
(d) $X_0$ is $\st A$ or $\st B$ with probability $\tfrac12$ each. Compute the distributions of $X_1$ and $X_2$.
(e) Given $X_1=\st B$, does the value of $X_0$ change the distribution of $X_2$? Which property says so?
\begin{solution}
(a) Entries non-negative, every row sums to $1$. Row $\st B$: after a trip starting in $\st B$, the next trip starts in $\st A$ with probability $0.2$, in $\st B$ with probability $0.6$, in $\st C$ with probability $0.2$.
(b) $p(\st A,\st C)\,p(\st C,\st B)=0.2\cdot0.3=0.06$.
(c) $(0.5,0.3,0.2)\PP=(0.25+0.06+0.06,\ 0.15+0.18+0.06,\ 0.10+0.06+0.08)=(0.37,0.39,0.24)$; $\Prob(X_2=\st C\mid X_0=\st A)=0.24=p(\st A,\st A)p(\st A,\st C)+p(\st A,\st B)p(\st B,\st C)+p(\st A,\st C)p(\st C,\st C)$: the sequences $\st{AAC},\st{ABC},\st{ACC}$.
(d) $\lambda=(\tfrac12,\tfrac12,0)$: $\lambda\PP=(0.35,0.45,0.20)$ and $\lambda\PP^2=(0.35,0.45,0.20)\PP=(0.325,0.435,0.24)$.
(e) No: $\Prob(X_2=j\mid X_1=\st B,X_0=i)=p(\st B,j)$ for every $i$. That is the Markov property.
% python3 sympy: P^2 row A = (37/100,39/100,6/25); lam P=(7/20,9/20,1/5); lam P^2=(13/40,87/200,6/25)
\end{solution}

\item The two-state chain $\PP=\begin{pmatrix}0.9&0.1\\0.3&0.7\end{pmatrix}$ on $\{0,1\}$ (the text source of B21).
(a) Find the eigenvalues of $\PP$ and the stationary distribution $\pi$.
(b) Show, by induction or from the lecture's formula for two-state chains, that
$\PP^n=\begin{pmatrix}\tfrac34&\tfrac14\\ \tfrac34&\tfrac14\end{pmatrix}+0.6^n\begin{pmatrix}\tfrac14&-\tfrac14\\-\tfrac34&\tfrac34\end{pmatrix}$, and check $n=1,2$.
(c) Compute $p^{(3)}(0,0)$ and $\Prob(X_3=0\mid X_0=1)$ exactly.
(d) Find the smallest $n$ for which every entry of $\PP^n$ is within $0.01$ of $\pi$.
\begin{solution}
(a) Trace $1.6$, so the eigenvalues are $1$ and $0.6$. $\pi\PP=\pi$: $0.1\pi_0=0.3\pi_1$, $\pi=(\tfrac34,\tfrac14)$.
(b) $n=1$: $\tfrac34+0.6\cdot\tfrac14=0.9$, $\tfrac14-0.6\cdot\tfrac14=0.1$, $\tfrac34-0.6\cdot\tfrac34=0.3$, $\tfrac14+0.6\cdot\tfrac34=0.7$. Induction: multiplying the first matrix by $\PP$ leaves it unchanged (its rows are $\pi$), and multiplying the second by $\PP$ multiplies it by the eigenvalue $0.6$ (its rows are proportional to $(1,-1)$, and $(1,-1)\PP=(0.6,-0.6)$). $n=2$: $\PP^2=\begin{pmatrix}0.84&0.16\\0.48&0.52\end{pmatrix}$.
(c) $p^{(3)}(0,0)=\tfrac34+\tfrac14\cdot0.216=0.804=\tfrac{201}{250}$; $p^{(3)}(1,0)=\tfrac34-\tfrac34\cdot0.216=0.588=\tfrac{147}{250}$.
(d) The largest deviation is $\tfrac34\cdot0.6^n$: $0.0126$ at $n=8$, $0.0076$ at $n=9$. So $n=9$.
% python3 sympy: P^2=[[21/25,4/25],[12/25,13/25]], P^3=[[201/250,49/250],[147/250,103/250]]; 0.75*0.6^8=0.01260, ^9=0.00756
\end{solution}

\item The vowel/consonant chain of both lectures: $S=\{\st V,\st C\}$, $\PP=\begin{pmatrix}0.2&0.8\\0.5&0.5\end{pmatrix}$, $\pi=(\tfrac5{13},\tfrac8{13})$.
(a) Compute $\PP^2$ and $\PP^3$.
(b) Compute the expected number of vowels among the next three letters, starting after a consonant and starting after a vowel, by linearity column by column; compare with the long-run value $3\pi_{\st V}$.
(c) Compute $\Prob(X_1=\st V,X_2=\st V\mid X_0=\st C)$ and compare with the stationary probability of the digraph $\st{VV}$ used in the Entropy lecture.
(d) Compute the expected number of letters until the first consonant, starting from $\st V$; then the expected number of letters until the pattern $\st{VC}$ is first completed, from $\st V$ and from $\st C$.
(e) Compute the mean return time to $\st C$ from $\pi$ and confirm it by a first-step argument.
\begin{solution}
(a) $\PP^2=\begin{pmatrix}0.44&0.56\\0.35&0.65\end{pmatrix}$, $\PP^3=\begin{pmatrix}0.368&0.632\\0.395&0.605\end{pmatrix}$.
(b) $\E_{\st C}\big[\sum_{n=1}^3\ind\{X_n=\st V\}\big]=\sum_{n=1}^3p^{(n)}(\st C,\st V)=0.5+0.35+0.395=1.245$; from $\st V$: $0.2+0.44+0.368=1.008$; long run $3\cdot\tfrac5{13}=1.154$, in between.
(c) $p(\st C,\st V)p(\st V,\st V)=0.5\cdot0.2=0.1$; stationary $\pi_{\st V}p(\st V,\st V)=\tfrac5{13}\cdot0.2=\tfrac1{13}=0.077$. A consonant just before makes a vowel, hence $\st{VV}$, likelier than average.
(d) From $\st V$ each letter is a consonant with probability $0.8$: geometric, mean $1/0.8=1.25$. Pattern $\st{VC}$: from $\st V$ the pattern completes at the first consonant, $k_{\st V}=1.25$; from $\st C$, $k_{\st C}=1+0.5k_{\st V}+0.5k_{\st C}$, so $k_{\st C}=2+k_{\st V}=3.25$.
(e) $1/\pi_{\st C}=\tfrac{13}8=1.625$. First step from $\st C$: with probability $0.5$ back at once, with probability $0.5$ at $\st V$ and then $1.25$ more letters: $1+0.5\cdot1.25=1.625$.
% python3: P^2=[[11/25,14/25],[7/20,13/20]], P^3=[[46/125,79/125],[79/200,121/200]]; 1.245, 1.008, 15/13; 1/10 vs 1/13; k_V=5/4, k_C=13/4; 13/8
\end{solution}

\subsection*{Stationary distributions \hfill\normalfont\small\color{soft}Lecture battery 8}
\item Random walks on graphs: $p(i,j)=1/d(i)$ for each neighbour $j$ of $i$.
(a) The path $0-1-2-3$. Find $\pi$ by detailed balance and verify the balance equation on the edges $(0,1)$ and $(1,2)$.
(b) The star with centre $0$ and leaves $1,2,3$. Find $\pi$; compute the mean return time to a leaf from $\pi$ and confirm it by a first-step argument.
(c) The Ehrenfest urn with $N=3$ balls. Find $\pi$ by detailed balance.
(d) The biased walk on the cycle $\Z_3$: $p(i,i+1)=\tfrac23$, $p(i,i-1)=\tfrac13$. Show that the uniform distribution is stationary but that detailed balance fails. What does the chain do that a reversible chain cannot?
\begin{solution}
(a) $\pi_i\propto d(i)$: $\pi=(1,2,2,1)/6$. Edge $(0,1)$: $\tfrac16\cdot1=\tfrac26\cdot\tfrac12$. Edge $(1,2)$: $\tfrac26\cdot\tfrac12=\tfrac26\cdot\tfrac12$.
(b) Degrees $(3,1,1,1)$: $\pi=(3,1,1,1)/6$; mean return time to leaf $1$ is $1/\pi_1=6$. First step: from the leaf to the centre ($1$ step); from the centre the walk picks leaf $1$ with probability $\tfrac13$, otherwise makes a wasted excursion of $2$ steps and is back at the centre; the number of wasted excursions is geometric with mean $2$; total $1+2\cdot2+1=6$.
(c) $\pi_i\,\frac{3-i}3=\pi_{i+1}\frac{i+1}3$, so $\pi_1=3\pi_0$, $\pi_2=\pi_1$, $\pi_3=\tfrac13\pi_2$: $\pi=(1,3,3,1)/8=\binom3i2^{-3}$.
(d) Each column of $\PP$ sums to $\tfrac23+\tfrac13=1$, so $\sum_i\tfrac13p(i,j)=\tfrac13$: uniform $\pi$ is stationary. Detailed balance: $\pi_0p(0,1)=\tfrac29\ne\tfrac19=\pi_1p(1,0)$. The chain has a net circulation ($\tfrac19$ per step anticlockwise around the cycle); run backwards it circulates the other way, so its law is not the same. A reversible chain has zero net flow on every edge.
% python3: path (1,2,2,1)/6; star (3,1,1,1)/6, return 6; Ehrenfest (1,3,3,1)/8; biased cycle 2/9 vs 1/9
\end{solution}

\item The taxi chain of C11.
(a) Solve $\pi\PP=\pi$, $\sum_i\pi_i=1$ exactly.
(b) $\PP^{10}$ has every row equal to $(0.3214,0.4286,0.2500)$ to four places. Confirm that this agrees with (a).
(c) Test detailed balance on the pair $(\st A,\st B)$. Is the chain reversible?
(d) The chain is irreducible and aperiodic. What do the convergence theorem and the ergodic theorem each say about $\pi_{\st C}=\tfrac14$?
\begin{solution}
(a) Columns $\st A$ and $\st C$: $0.5\pi_{\st A}+0.2\pi_{\st B}+0.3\pi_{\st C}=\pi_{\st A}$ and $0.2\pi_{\st A}+0.2\pi_{\st B}+0.4\pi_{\st C}=\pi_{\st C}$, i.e.\ $5\pi_{\st A}=2\pi_{\st B}+3\pi_{\st C}$ and $3\pi_{\st C}=\pi_{\st A}+\pi_{\st B}$. With $\pi_{\st A}+\pi_{\st B}+\pi_{\st C}=1$ the second gives $\pi_{\st C}=\tfrac14$, then $5\pi_{\st A}=2\pi_{\st B}+\tfrac34$ and $\pi_{\st A}+\pi_{\st B}=\tfrac34$ give $\pi=(\tfrac9{28},\tfrac37,\tfrac14)=(0.3214,0.4286,0.25)$.
(b) $\tfrac9{28}=0.32143$, $\tfrac37=0.42857$: the rows of $\PP^{10}$ are $\pi$ to four places, as the convergence theorem predicts.
(c) $\pi_{\st A}p(\st A,\st B)=\tfrac9{28}\cdot0.3=\tfrac{27}{280}$ but $\pi_{\st B}p(\st B,\st A)=\tfrac37\cdot0.2=\tfrac{24}{280}$. Not reversible (unlike the weather chain of the lecture): more trips go $\st A\to\st B$ than $\st B\to\st A$ in the long run.
(d) Convergence: $\Prob(X_n=\st C)\to\tfrac14$ whatever the starting zone. Ergodic: along a single sequence of trips, the fraction starting in $\st C$ tends to $\tfrac14$ with probability one, and consecutive $\st C$-starts are $1/\pi_{\st C}=4$ trips apart on average.
% python3 sympy: pi=(9/28,3/7,1/4); P^10 rows (0.32143,0.42857,0.25); 27/280 vs 24/280
\end{solution}

\item PageRank with damping $d=0.85$.
(a) Three pages: $\st A\to\st B$, $\st A\to\st C$, $\st B\to\st C$, $\st C\to\st A$. Write $\mathbf H$ and $\PP=0.85\,\mathbf H+0.05\,\mathbf J$.
(b) Solve $\pi\PP=\pi$ and rank the pages. $\st A$ and $\st B$ each have exactly one in-link; why does $\st A$ rank higher?
(c) Find $\pi$ without damping ($d=1$).
(d) Add a page $\st D$ with links $\st C\to\st D$ (so $\st C$ now links to $\st A$ and $\st D$) and $\st D\to\st D$ only. Find $\pi$ without damping, and with $d=0.85$. What does $\st D$ do to the surfer, and how much does damping help?
\begin{solution}
(a) $\mathbf H=\begin{pmatrix}0&\tfrac12&\tfrac12\\0&0&1\\1&0&0\end{pmatrix}$, $\PP=\begin{pmatrix}0.05&0.475&0.475\\0.05&0.05&0.9\\0.9&0.05&0.05\end{pmatrix}$.
(b) $\pi\PP=\pi$ with $\sum\pi=1$: $\pi_{\st A}=0.05+0.85\pi_{\st C}$, $\pi_{\st B}=0.05+0.425\pi_{\st A}$, $\pi_{\st C}=0.05+0.425\pi_{\st A}+0.85\pi_{\st B}$; solving, $\pi=(\tfrac{686}{1769},\tfrac{380}{1769},\tfrac{703}{1769})=(0.388,0.215,0.397)$: $\st C>\st A>\st B$. $\st A$'s one in-link is the \emph{only} link of the top page $\st C$, so it carries all of $\st C$'s weight; $\st B$'s in-link is one of two from $\st A$ and carries half of $\st A$'s.
(c) $\pi\mathbf H=\pi$: $\pi_{\st A}=\pi_{\st C}$, $\pi_{\st B}=\tfrac12\pi_{\st A}$, $\pi_{\st C}=\tfrac12\pi_{\st A}+\pi_{\st B}$: $\pi=(\tfrac25,\tfrac15,\tfrac25)$.
(d) Without damping, $\st D$ is absorbing: once the surfer arrives (which happens, since $\st C$ is visited again and again) it never leaves, and the only stationary distribution is $\pi=(0,0,0,1)$. With $d=0.85$, solving $\pi\PP=\pi$ gives $\pi=(0.101,0.080,0.148,0.671)$: $\st D$ still holds two thirds of the mass, because each visit lasts a geometric number of steps with mean $1/0.15=6.7$ and every jump lands on $\st D$ a quarter of the time. Damping makes the problem finite but does not neutralise a self-linking trap; real search engines treat such pages separately.
% python3 sympy: 3-page pi = (686,380,703)/1769; undamped (2/5,1/5,2/5); 4-page trap: undamped (0,0,0,1), damped (0.1006,0.0803,0.1485,0.6707)
\end{solution}

\subsection*{Classes, periods, recurrence \hfill\normalfont\small\color{soft}Lecture battery 9}
\item A chain on $S=\{1,\dots,6\}$ with
\[
  \PP=\begin{pmatrix}0&1&0&0&0&0\\1&0&0&0&0&0\\0.2&0&0.4&0&0.4&0\\0&0&0.5&0.5&0&0\\0&0&0&0&0&1\\0&0&0&0&0.5&0.5\end{pmatrix}.
\]
(a) Find the communicating classes, say which are closed, give the period of each class, and classify each state as recurrent or transient.
(b) Compute the probability that the chain started at $3$, and at $4$, ends up in the class $\{1,2\}$.
(c) Compute the expected time to leave $\{3,4\}$ from $3$ and from $4$.
(d) Find all stationary distributions.
(e) Started at $3$, does the distribution of $X_n$ converge? Describe its behaviour for large $n$.
\begin{solution}
(a) $\{1,2\}$: closed, $1\to2\to1$ only, period $2$. $\{5,6\}$: closed, $p(6,6)>0$, period $1$. $\{3\}$: $3\to1$ and $3\to5$ leave it, so not closed; $p(3,3)>0$, period $1$. $\{4\}$: $4\to3$ leaves it, not closed, period $1$. States $1,2,5,6$ recurrent (finite closed classes); $3,4$ transient.
(b) $h_3=0.2\cdot1+0.4h_3+0.4\cdot0$, so $h_3=\tfrac13$; $h_4=0.5h_3+0.5h_4$, so $h_4=h_3=\tfrac13$. With probability $\tfrac23$ the chain ends in $\{5,6\}$.
(c) $k_3=1+0.4k_3$: $k_3=\tfrac53$; $k_4=1+0.5k_3+0.5k_4$: $k_4=2+k_3=\tfrac{11}3$.
(d) On $\{1,2\}$ the unique stationary distribution is $(\tfrac12,\tfrac12)$; on $\{5,6\}$, $\pi_5=0.5\pi_6$ gives $(\tfrac13,\tfrac23)$. Transient states carry no mass: if $\pi_3+\pi_4>0$ then, since $\{3,4\}$ is left for good with positive probability at every step and never re-entered from $\{1,2\}$ or $\{5,6\}$, $\pi\PP^n$ puts mass $\to0$ on $\{3,4\}$, contradicting $\pi\PP^n=\pi$; so $\pi_3=\pi_4=0$. All stationary distributions: $\alpha(\tfrac12,\tfrac12,0,0,0,0)+(1-\alpha)(0,0,0,0,\tfrac13,\tfrac23)$, $\alpha\in[0,1]$.
(e) No. The mass $\tfrac13$ that enters $\{1,2\}$ alternates between the two states forever (period $2$), so $\Prob_3(X_n=1)$ oscillates; the mass $\tfrac23$ in $\{5,6\}$ settles to $(\tfrac29,\tfrac49)$. Both hypotheses of the convergence theorem fail: the chain is reducible, and one closed class is periodic.
% python3 sympy: h3=h4=1/3; k3=5/3, k4=11/3; pi on {5,6}=(1/3,2/3)
\end{solution}

\item Walks on the cycle $\Z_L=\{0,\dots,L-1\}$, addition modulo $L$.
(a) $L=6$, steps $\pm1$ with probability $\tfrac12$ each. Compute $p^{(n)}(0,0)$ for $n=1,\dots,4$, give the period, and describe the rows of $\PP^n$ for large $n$ (is $\pi$ stationary; do the rows converge to it?).
(b) $L=7$, steps $\pm1$. Give the period and justify it with two round-trip lengths.
(c) $L=6$, steps $+1$ or $+2$ with probability $\tfrac12$ each. Is the chain irreducible? Aperiodic? What is $\pi$, and do the rows of $\PP^n$ converge?
(d) $L=6$, steps $+2$ or $+4$ with probability $\tfrac12$ each. Find the communicating classes and the period within each.
(e) $L=4$, step $+1$ with probability one. Find $\pi$, the period, and say which of the two long-run theorems applies.
\begin{solution}
(a) $p^{(n)}(0,0)=0,\tfrac12,0,\tfrac38$ ($n=4$: four steps with two of each sign, $\binom42/16$; a net displacement of $\pm6$ is impossible in four steps): every return takes an even number of steps, period $2$. The uniform $\pi=\tfrac16(1,\dots,1)$ is stationary (each column receives $\tfrac12$ from two states), but the rows of $\PP^n$ do not converge to it: from $0$ the walk is on $\{0,2,4\}$ at even times and on $\{1,3,5\}$ at odd times, and row $0$ of $\PP^n$ alternates between (approximately) $(\tfrac13,0,\tfrac13,0,\tfrac13,0)$ and $(0,\tfrac13,0,\tfrac13,0,\tfrac13)$.
(b) Round trips of length $2$ ($+1,-1$) and $7$ (seven steps $+1$): $\gcd(2,7)=1$, period $1$; here the rows of $\PP^n$ do converge to the uniform $\pi$.
% python3 numpy: L=6 p^(n)(0,0)=0,1/2,0,3/8,0,11/32,...; row 0 of P^100 = (1/3,0,1/3,0,1/3,0), P^101 = (0,1/3,0,1/3,0,1/3); L=7 P^200 rows uniform 1/7
(c) Irreducible: $+1$ steps reach every state. Round trips of length $3$ ($+2,+2,+2$) and $4$ ($+1,+1,+2,+2$): aperiodic. Every column of $\PP$ receives $\tfrac12$ from two states, so the uniform $\pi=\tfrac16(1,\dots,1)$ is stationary, and by the convergence theorem every row of $\PP^n$ tends to it.
(d) From $0$ only $0,2,4$ are reachable: two closed classes, $\{0,2,4\}$ and $\{1,3,5\}$; the chain is reducible. Within a class, round trips of length $2$ ($+2,+4$) and $3$ ($+2,+2,+2$): period $1$.
(e) Deterministic rotation: $\pi$ uniform (the only solution of $\pi_{i+1}=\pi_i$), period $4$. The ergodic theorem applies (irreducible): each state is visited exactly a quarter of the time. The convergence theorem does not: $\Prob_0(X_n=0)$ is $1$ or $0$ according to $n\bmod4$.
% python3 numpy: L=6 (+1,+2) returns 3,4,...; L=6 (+2,+4) reachable {0,2,4}, returns 2,3
\end{solution}

\subsection*{Hitting probabilities and times \hfill\normalfont\small\color{soft}Lecture battery 10}
\item Tennis at deuce. From deuce the server wins each point with probability $p$, independently. States: $\st D$ (deuce), $\st I$ (advantage server), $\st O$ (advantage receiver), $\st W$ (game won), $\st L$ (game lost). From $\st D$ the next state is $\st I$ or $\st O$; from $\st I$ it is $\st W$ or $\st D$; from $\st O$ it is $\st D$ or $\st L$.
(a) Write the transition matrix.
(b) Write the hitting system for $h_i=\Prob_i(\text{reach }\st W)$ and solve it; show $h_{\st D}=\dfrac{p^2}{p^2+q^2}$.
(c) Write the system for $k_i=\E_i[\text{points until the game ends}]$ and show $k_{\st D}=\dfrac{2}{1-2pq}$.
(d) Evaluate $h_{\st D}$ and $k_{\st D}$ for $p=0.5,0.55,0.6,0.7$. What does a $55/45$ edge per point become per game?
(e) What is the head-on computation of $h_{\st D}$, and why does the first-step system avoid it?
\begin{solution}
(a) Order $\st D,\st I,\st O,\st W,\st L$:
$\PP=\begin{pmatrix}0&p&q&0&0\\q&0&0&p&0\\p&0&0&0&q\\0&0&0&1&0\\0&0&0&0&1\end{pmatrix}$.
(b) $h_{\st W}=1$, $h_{\st L}=0$, $h_{\st D}=ph_{\st I}+qh_{\st O}$, $h_{\st I}=p+qh_{\st D}$, $h_{\st O}=ph_{\st D}$. Substituting, $h_{\st D}=p^2+pq\,h_{\st D}+pq\,h_{\st D}$, so $h_{\st D}(1-2pq)=p^2$ and, as $1-2pq=p^2+q^2$, $h_{\st D}=\frac{p^2}{p^2+q^2}$. For $p=0.6$: $h_{\st D}=\tfrac9{13}$, $h_{\st I}=\tfrac{57}{65}$, $h_{\st O}=\tfrac{27}{65}$.
(c) $k_{\st W}=k_{\st L}=0$, $k_{\st D}=1+pk_{\st I}+qk_{\st O}$, $k_{\st I}=1+qk_{\st D}$, $k_{\st O}=1+pk_{\st D}$. Then $k_{\st D}=1+p+q+2pq\,k_{\st D}=2+2pq\,k_{\st D}$, so $k_{\st D}=\frac2{1-2pq}$; for $p=0.6$, $k_{\st D}=\tfrac{50}{13}=3.85$ points.
(d) $h_{\st D}=0.5,\ 0.599,\ 0.692,\ 0.845$ and $k_{\st D}=4,\ 3.96,\ 3.85,\ 3.45$. A $55/45$ point becomes a $60/40$ game from deuce; the edge compounds because the game only ends when one player is two points ahead.
(e) The game ends after $2m$ more points with probability $(2pq)^{m-1}(p^2+q^2)$, and the server wins with probability $\sum_{m\ge1}(2pq)^{m-1}p^2$: an infinite series. The first-step system uses the restart at deuce to write $h_{\st D}$ in terms of itself.
% python3 sympy: p=0.6: h=(9/13,57/65,27/65), k=(50/13,33/13,43/13); table 0.5/4, 0.5990/3.9604, 0.6923/3.8462, 0.8448/3.4483
\end{solution}

\item The taxi chain of C11 again.
(a) Compute the expected number of trips until the next trip starting in $\st C$, from $\st A$ and from $\st B$.
(b) Compute the mean return time to $\st C$ by first-step analysis and compare with $1/\pi_{\st C}$.
(c) Make $\st C$ absorbing and compute the fundamental matrix $\mathbf N$; recover (a) from $\mathbf N\mathbf 1$ and read off the expected number of trips starting in $\st B$ before the next $\st C$-start, when the current trip starts in $\st A$.
\begin{solution}
(a) $k_{\st C}=0$; $k_{\st A}=1+0.5k_{\st A}+0.3k_{\st B}$, $k_{\st B}=1+0.2k_{\st A}+0.6k_{\st B}$. The second gives $0.4k_{\st B}=1+0.2k_{\st A}$, the first $0.5k_{\st A}=1+0.3k_{\st B}$; solving, $k_{\st A}=k_{\st B}=5$ trips.
(b) $\E_{\st C}[T_{\st C}]=1+0.3k_{\st A}+0.3k_{\st B}+0.4\cdot0=1+1.5+1.5=4=1/\pi_{\st C}$.
(c) $\mathbf Q=\begin{pmatrix}0.5&0.3\\0.2&0.6\end{pmatrix}$, $\mathbf I-\mathbf Q=\begin{pmatrix}0.5&-0.3\\-0.2&0.4\end{pmatrix}$, determinant $0.14$, $\mathbf N=\frac1{0.14}\begin{pmatrix}0.4&0.3\\0.2&0.5\end{pmatrix}=\begin{pmatrix}\tfrac{20}7&\tfrac{15}7\\ \tfrac{10}7&\tfrac{25}7\end{pmatrix}$; $\mathbf N\mathbf 1=(5,5)$. $\mathbf N_{\st A\st B}=\tfrac{15}7=2.14$ trips start in $\st B$, on average, before the next $\st C$-start.
% python3 sympy: kA=kB=5; return 4; N=[[20/7,15/7],[10/7,25/7]]
\end{solution}

\subsection*{Convergence, the ergodic theorem, absorbing chains \hfill\normalfont\small\color{soft}Lecture batteries 11--12}
\item Long-run behaviour of the taxi chain of C11 and C15.
(a) The eigenvalues of $\PP$ are $1$, $0.3$ and $0.2$ (check the trace and the determinant). How fast do the rows of $\PP^n$ approach $\pi$? The maximal deviations are $0.179,0.051,0.015$ at $n=1,2,3$; compare with $0.3^n$.
(b) State what fraction of trips start in each zone in the long run, and which theorem says so.
(c) What fraction of trips go from $\st A$ to $\st B$ in the long run?
(d) Trips starting in $\st A,\st B,\st C$ earn $\pounds10,\pounds8,\pounds12$. What is the long-run average fare per trip?
(e) For the rotation on $\Z_4$ of C18(e), the long-run fraction of time at each state is $\tfrac14$ but $\Prob_0(X_n=0)$ does not converge. Which hypothesis of the convergence theorem fails, and why does the ergodic theorem not need it?
\begin{solution}
(a) Trace $0.5+0.6+0.4=1.5=1+0.3+0.2$; determinant $0.5(0.24-0.06)-0.3(0.08-0.06)+0.2(0.06-0.18)=0.06=1\cdot0.3\cdot0.2$. The error shrinks by a factor of about $0.3$ per step: $0.179,0.051,0.015$ have successive ratios $0.29,0.30$.
(b) $\pi=(\tfrac9{28},\tfrac37,\tfrac14)$: $32.1\%$, $42.9\%$, $25\%$ of trips, with probability one along the actual sequence of trips (ergodic theorem, irreducible finite chain).
(c) Ergodic theorem for the pair chain $(X_n,X_{n+1})$: $\pi_{\st A}p(\st A,\st B)=\tfrac9{28}\cdot0.3=\tfrac{27}{280}=9.6\%$.
(d) $\sum_i\pi_if(i)=\tfrac9{28}\cdot10+\tfrac37\cdot8+\tfrac14\cdot12=\tfrac{135}{14}=\pounds9.64$ per trip.
(e) Aperiodicity fails (period $4$). The ergodic theorem counts visits along a path; the fraction of time at $0$ is exactly $\tfrac14$ whether or not the visits come at regular intervals. The convergence theorem asks where the chain is at one fixed late time, and a periodic chain still remembers the starting time.
% python3 numpy: eigenvalues 1, 0.3, 0.2; max|P^n-pi| = 0.1786, 0.0514, 0.0154; 27/280; 135/14=9.643
\end{solution}

\item The fundamental matrix.
(a) Fair ruin chain on $\{0,\dots,5\}$, transient states $1,2,3,4$. Write $\mathbf Q$ and $\mathbf R$, compute $\mathbf N=(\mathbf I-\mathbf Q)^{-1}$, and check $\mathbf N\mathbf 1$ against $i(5-i)$ and the target column of $\mathbf N\mathbf R$ against $i/5$. Starting from $\pounds1$, how many times on average is the fortune $\pounds3$ before the game ends, and how many times $\pounds1$ (counting the start)?
(b) Tennis, $p=0.6$ (C19). Compute $\mathbf N$ for the transient states $\st D,\st I,\st O$ and confirm $h_{\st D}=\tfrac9{13}$ and $k_{\st D}=\tfrac{50}{13}$ from $\mathbf N\mathbf R$ and $\mathbf N\mathbf 1$. Interpret $\mathbf N_{\st D\st D}$ and check it against the visits criterion $\mathbf N_{ii}=1/(1-f_i)$ with $f_{\st D}=\Prob_{\st D}(\text{return to }\st D)$.
\begin{solution}
(a) $\mathbf Q$ is $4\times4$ with $\tfrac12$ on the two off-diagonals and zeros elsewhere; $\mathbf R$ has $\tfrac12$ in row $1$, column $0$ and in row $4$, column $5$.
$\mathbf N=\tfrac15\begin{pmatrix}8&6&4&2\\6&12&8&4\\4&8&12&6\\2&4&6&8\end{pmatrix}$ (check: $(\mathbf I-\mathbf Q)\mathbf N=\mathbf I$ row by row). $\mathbf N\mathbf 1=(4,6,6,4)=i(5-i)$; $\mathbf N\mathbf R$ has target column $(\tfrac15,\tfrac25,\tfrac35,\tfrac45)=i/5$. From $\pounds1$: $\mathbf N_{13}=\tfrac45$ visits to $\pounds3$, $\mathbf N_{11}=\tfrac85$ visits to $\pounds1$.
(b) $\mathbf Q=\begin{pmatrix}0&0.6&0.4\\0.4&0&0\\0.6&0&0\end{pmatrix}$, $\mathbf R=\begin{pmatrix}0&0\\0.6&0\\0&0.4\end{pmatrix}$ (columns $\st W,\st L$), $\mathbf N=\tfrac1{13}\begin{pmatrix}25&15&10\\10&19&4\\15&9&19\end{pmatrix}$; $\mathbf N\mathbf 1=(\tfrac{50}{13},\tfrac{33}{13},\tfrac{43}{13})$ and $\mathbf N\mathbf R$ has first column $(\tfrac9{13},\tfrac{57}{65},\tfrac{27}{65})$, as in C19. $\mathbf N_{\st D\st D}=\tfrac{25}{13}=1.92$ is the expected number of deuces in the game, the first included; $f_{\st D}=2pq=0.48$ (one point each way), and $1/(1-0.48)=\tfrac{25}{13}$.
% python3 sympy: N ruin N=5 = (1/5)[[8,6,4,2],[6,12,8,4],[4,8,12,6],[2,4,6,8]], N1=(4,6,6,4), NR target (1,2,3,4)/5; tennis N=(1/13)[[25,15,10],[10,19,4],[15,9,19]]
\end{solution}
\end{enumerate}
